% ======================================================================
% PEC VOLUNTARIA (UNED) — Psicología del Pensamiento
% Resolución de problemas por insight: réplica adaptada de Knoblich et al. (1999)
% Máx. 1,5 páginas (sin contar portada y referencias)
% ======================================================================
\documentclass[10pt,a4paper]{article}

\usepackage[utf8]{inputenc}
\usepackage[spanish]{babel}
\usepackage{csquotes}
\usepackage{geometry}
\usepackage{microtype}
\usepackage{graphicx}
\usepackage{booktabs}
\usepackage{caption}
\usepackage{subcaption}
\usepackage{enumitem}
\usepackage{hyperref}
\usepackage{float}

\usepackage[backend=biber,style=apa,sortcites=true]{biblatex}
\addbibresource{referencias.bib}

\geometry{margin=1.8cm}
\setlength{\parindent}{0pt}
\setlength{\parskip}{0.25em}
\captionsetup{font=small,skip=2pt}
\hypersetup{colorlinks=true,linkcolor=black,urlcolor=blue,citecolor=black}

% Figuras generadas por scripts (rutas relativas desde /informe)
\graphicspath{{../data/gr\'aficos/}}

\title{\vspace{-1.2em}\textbf{Resolución de problemas por insight:}\\Réplica experimental adaptada basada en Knoblich et al. (1999)}
\author{\textbf{Estudiante:} [Tu nombre] \quad \textbf{Asignatura:} Psicología del Pensamiento (UNED) \\ \textbf{Fecha:} 03/05/2026}
\date{}

\begin{document}
\maketitle
\vspace{-1.2em}

\section{Objetivos e hipótesis}
\textbf{Objetivo.} Replicar de forma adaptada el Experimento 1A de \textcite{knoblich1999} usando problemas aritméticos de cerillas con numeración romana, comparando (a) la dificultad inicial entre tipos de problema (A--D) y (b) posibles efectos de transferencia tras una única experiencia previa (comparación Bloque 1 vs. Bloque 2).

\textbf{Hipótesis (descriptivas).}
\begin{itemize}[leftmargin=*, itemsep=0.1em]
    \item H1: En el Bloque 1, la tasa de solución y/o el tiempo medio difieren entre tipos (A--D), reflejando distintos requisitos de reestructuración (relajación de restricciones / descomposición de ``chunks'').
    \item H2: En el Bloque 2, se observa transferencia: mayor tasa de solución y/o menor tiempo medio respecto al Bloque 1, especialmente en tipos con reestructuración similar.
\end{itemize}

\section{Método}
\textbf{Diseño.} Intra-sujeto 2 (\textit{Bloque}: 1, 2) $\times$ 4 (\textit{Tipo de problema}: A, B, C, D). Variables dependientes: (a) \textit{resolución} (S/N) y (b) \textit{tiempo de resolución} (segundos; resumido en minutos).

\textbf{Participantes.} $N=10$ adultos (códigos anónimos P01--P10), reclutados por conveniencia. Se verificó familiaridad mínima con números romanos (I--XX). No se recogieron nombres reales.

\textbf{Materiales.} Plataforma web con representación de cerillas y numeración romana (app). Se utilizaron exactamente las ecuaciones requeridas en la PEC (Tabla 1 de \textcite{knoblich1999}), 4 en Bloque 1 y 4 en Bloque 2. Regla: mover \textit{una única cerilla} para convertir una ecuación falsa en verdadera. Tiempo máximo: 5 min por problema.

\textbf{Procedimiento.} El experimentador supervisó la tarea (presencial o videollamada), controló el tiempo y registró: orden de presentación (por bloque), acierto (S/N), tiempo (s) y observaciones cualitativas. El orden dentro de cada bloque varió entre participantes (9 secuencias distintas en cada bloque). Los datos brutos se consolidaron en la plantilla Excel (anexo obligatorio).

\section{Resultados}
La Tabla~\ref{tab:resumen} resume el porcentaje de problemas resueltos y el tiempo medio (min) por tipo y bloque. La frecuencia acumulada de soluciones (Fig.~\ref{fig:freq}) se calculó contabilizando \textit{solo problemas resueltos} (frecuencia acumulada de soluciones), y el tiempo medio y tasa de solución se muestran en la Fig.~\ref{fig:tiempo}.

\begin{table}[t]
\centering
\caption{Resumen descriptivo por bloque y tipo (N=10)}
\label{tab:resumen}
\vspace{-0.2em}
\small
\begin{tabular}{cccccc}
\toprule
\textbf{Bloque} & \textbf{Tipo} & \textbf{Resueltos} & \textbf{Total} & \textbf{Tasa (\%)} & \textbf{Tiempo medio (min)} \\
\midrule
1 & A & 8 & 10 & 80.0 & 1.49 \\
1 & B & 6 & 10 & 60.0 & 2.49 \\
1 & C & 5 & 10 & 50.0 & 2.64 \\
1 & D & 9 & 10 & 90.0 & 1.45 \\
\midrule
2 & A & 8 & 10 & 80.0 & 1.15 \\
2 & B & 4 & 10 & 40.0 & 2.07 \\
2 & C & 5 & 10 & 50.0 & 2.79 \\
2 & D & 9 & 10 & 90.0 & 1.73 \\
\bottomrule
\end{tabular}
\vspace{-0.6em}
\end{table}

\begin{figure}[t]
\centering
\begin{subfigure}[t]{0.49\linewidth}
    \centering
    \includegraphics[width=\linewidth,height=4.0cm,keepaspectratio]{fig_bloque_1_frecuencia.png}
    \caption{Bloque 1}
\end{subfigure}
\begin{subfigure}[t]{0.49\linewidth}
    \centering
    \includegraphics[width=\linewidth,height=4.0cm,keepaspectratio]{fig_bloque_2_frecuencia.png}
    \caption{Bloque 2}
\end{subfigure}
\caption{Frecuencia acumulada de soluciones por tipo de problema (A--D) y bloque.}
\label{fig:freq}
\vspace{-0.8em}
\end{figure}

\begin{figure}[t]
\centering
\begin{subfigure}[t]{0.49\linewidth}
    \centering
    \includegraphics[width=\linewidth,height=4.0cm,keepaspectratio]{fig_bloque_1_tiempo_y_tasa.png}
    \caption{Bloque 1}
\end{subfigure}
\begin{subfigure}[t]{0.49\linewidth}
    \centering
    \includegraphics[width=\linewidth,height=4.0cm,keepaspectratio]{fig_bloque_2_tiempo_y_tasa.png}
    \caption{Bloque 2}
\end{subfigure}
\caption{Tiempo medio de resolución (min) y tasa de solución (\%) por tipo y bloque.}
\label{fig:tiempo}
\vspace{-1.0em}
\end{figure}

\section{Discusión}
El patrón general observado muestra que el Tipo D fue el más fácil en ambos bloques (90\% de solución), mientras que los tipos B/C tendieron a mostrar menor rendimiento (especialmente B en Bloque 2). En términos de transferencia, los datos no muestran una mejora uniforme del Bloque 1 al Bloque 2 (p.ej., A mantiene 80\%, D mantiene 90\%, B disminuye de 60\% a 40\%), aunque algunos tiempos medios cambian entre bloques. En línea con el marco de \textcite{knoblich1999}, estas diferencias pueden interpretarse como resultado de restricciones implícitas y de la necesidad de reestructuración (relajación de restricciones y/o descomposición de unidades perceptivas) que varía por tipo de problema.

Como limitaciones, el tamaño muestral es reducido y cada condición contiene un único ítem por tipo y bloque (diseño adaptado), por lo que el objetivo es descriptivo: evaluar si emerge un patrón general compatible con el descrito por \textcite{knoblich1999}. Futuras réplicas podrían aumentar el número de participantes y añadir más ítems por condición para estabilizar las estimaciones.

\vspace{0.2em}
\textbf{Anexo obligatorio (datos brutos):} plantilla Excel cumplimentada con los 8 problemas por participante.

\printbibliography
\end{document}